\section{Atores Primitivos}

As funcionalidades mencionadas são compostas de algumas capacidades primitivas:

\begin{itemize}
    \item Operações do sistema de arquivos
    \item Requisições HTTP
    \item Manipulação de variáveis de ambiente
    \item Controle de terminal
    \item Integração com shell
\end{itemize}

Capacidades primitivas podem ser entendidas como blocos de construção do sistema que não fazem parte da lógica de negócio e dependem principalmente de interagir com sistemas externos (rede, shell, arquivos, etc). Cada uma dessas capacidades pode ser traduzida para um ator em nossa aplicação seguindo os princípios SOLID, especialmente:

\begin{itemize}
    \item \textbf{Responsabilidade Única}: cada ator lida com um e apenas um domínio
    \item \textbf{Segregação de Interface}: ao dividir em diferentes atores, a injeção de dependência é mais granular, então clientes dependem apenas de interfaces que importam para eles
\end{itemize}

Esses atores serão chamados de atores primitivos, já que não dependem de outros atores.
